\DocumentMetadata{
  pdfversion=2.0,
  pdfstandard=ua-2,
  lang=en-US,
  tagging=on,
  tagging-setup={math/setup=mathml-SE,tabsorder=structure}
}
\documentclass[11pt,letterpaper]{article}
\usepackage[margin=0.68in,includefoot]{geometry}
\usepackage{newcomputermodern}
\usepackage{amsmath,amscd,mathtools}
\providecommand{\square}{\mdlgwhtsquare}
\usepackage[hidelinks]{hyperref}
\hypersetup{
  pdftitle={Numerical Analysis First-Year Exam, January 2026},
  pdfauthor={University of Florida Department of Mathematics},
  pdfsubject={Graduate mathematics examination},
  pdfdisplaydoctitle=true,
  bookmarksopen=true
}
\setcounter{secnumdepth}{0}
\setlength{\parindent}{0pt}
\setlength{\parskip}{0.28em}
\setlength{\emergencystretch}{3em}
\setlength{\leftmargini}{2.15em}
\setlength{\leftmarginii}{2.65em}
\setlength{\leftmarginiii}{3em}
\renewcommand{\labelenumi}{\textbf{\theenumi.}}
\pagestyle{empty}
\raggedbottom
\allowdisplaybreaks
\newenvironment{examproblems}[1]{%
  \begin{enumerate}%
  \setcounter{enumi}{#1}%
  \setlength{\topsep}{0.35em}%
  \setlength{\partopsep}{0pt}%
  \setlength{\itemsep}{0.48em}%
  \setlength{\parsep}{0.18em}%
}{\end{enumerate}}
\newenvironment{examsubparts}{%
  \begin{enumerate}%
  \setlength{\topsep}{0.18em}%
  \setlength{\partopsep}{0pt}%
  \setlength{\itemsep}{0.16em}%
  \setlength{\parsep}{0.1em}%
}{\end{enumerate}}
\newenvironment{examinstructions}{%
  \par\begingroup\itshape
}{%
  \par\endgroup\vspace{0.18em}
}
\makeatletter
\renewcommand\section{\@startsection{section}{1}{0pt}%
  {-1.25ex plus -.25ex minus -.1ex}%
  {0.55ex plus .1ex}%
  {\normalfont\large\bfseries}}
\renewcommand{\@maketitle}{%
  \newpage\null\vskip -1.2em%
  {\footnotesize\bfseries
    \href{https://www.ufl.edu/}{University of Florida}, \href{https://math.ufl.edu/}{Department of Mathematics}\par}%
  \vskip 0.18em%
  \tagstructbegin{tag=Title}%
    {\LARGE\bfseries\href{https://gma.math.ufl.edu/exams/numerical-analysis-first-year/}{Numerical Analysis First-Year Exam}, January 2026\par}%
  \tagstructend%
  \vskip 0.35em%
  \tagmcbegin{artifact}%
    \hrule height 1pt%
  \tagmcend%
  \vskip 0.55em%
}
\makeatother
\title{Numerical Analysis First-Year Exam, January 2026}
\author{}
\date{}
\begin{document}
\maketitle
\thispagestyle{empty}
\begin{examinstructions}
Do 4 (four) problems.
\end{examinstructions}
\begin{examproblems}{0}
\item
Prove each of the following.
\begin{examsubparts}
\item[\textnormal{(a)}]
The matrix 2-norm is invariant under unitary transformation.
\item[\textnormal{(b)}]
If \(n\times n\) matrix~\(A\) is diagonal, then for any matrix~\(p\)-norm \(\lVert\cdot\rVert_p\), it holds that \(\lVert A\rVert_p=\max_{1\le i\le n}|a_{ii}|\).
\end{examsubparts}
\item
Suppose \(A\in\mathbb{C}^{m\times n}\). Prove or provide a counterexample for each of the following.
\begin{examsubparts}
\item[\textnormal{(a)}]
The Frobenius norm of a matrix is induced by (subordinate to) a vector norm.
\item[\textnormal{(b)}]
\(\lVert A\rVert_2\le\lVert A\rVert_F\le\sqrt{n}\lVert A\rVert_2\).
\end{examsubparts}
\item
Let \(A\in\mathbb{C}^{n\times n}\). You may use the fact (if you find it useful) That~\(A\) has a Schur decomposition. Prove or provide a counterexample for each of the following.
\begin{examsubparts}
\item[\textnormal{(a)}]
If~\(A\) is both normal and triangular then it is diagonal.
\item[\textnormal{(b)}]
If~\(A\) is normal then the 2-norm of~\(A\) is equal to it’s spectral radius \(\rho(A)\).
\end{examsubparts}
\item
This problem has two parts.
\begin{examsubparts}
\item[\textnormal{(a)}]
Let \(r_i=\sum_{j=1,j\ne i}^{n}|a_{ij}|\). Let \(D_i\) be the disk in~\(\mathbb{C}\) with center \(a_{ii}\) and radius \(r_i\). Prove that if~\(\lambda\) is an eigenvalue of~\(A\), then \(\lambda\in\bigcup_i D_i\); in other words, \(\lambda\) is in at least one of the disks \(D_i\).
\item[\textnormal{(b)}]
Determine all eigenvalues of the Householder reflector \(H(w)\), including their multiplicities. Justify your answer.
\end{examsubparts}
\item
Define the matrices~\(A\) and~\(B\) by 
\[
A=\begin{pmatrix}1&3&0\\1&3&0\\1&3&0\\1&3&0\end{pmatrix},\qquad B=\begin{pmatrix}1&0&1\\0&2&0\\1&0&0\\0&0&1\end{pmatrix}.
\]
\begin{examsubparts}
\item[\textnormal{(a)}]
Find a full singular value decomposition of~\(A\).
\item[\textnormal{(b)}]
Find an economy QR decomposition of~\(B\).
\end{examsubparts}
\end{examproblems}
\end{document}
